\documentclass[10pt]{beamer}
\usetheme{metropolis}
\usepackage{graphicx}
\usepackage{listings}
\usepackage{booktabs}
\usepackage{xcolor}
\usepackage{hyperref}

\definecolor{codebg}{HTML}{F5F5F5}
\definecolor{codegreen}{HTML}{2E7D32}
\definecolor{codegray}{HTML}{757575}
\definecolor{codeblue}{HTML}{1565C0}

\lstdefinestyle{rstyle}{
  language=R, backgroundcolor=\color{codebg}, basicstyle=\ttfamily\scriptsize,
  keywordstyle=\color{codeblue}\bfseries, stringstyle=\color{codegreen},
  commentstyle=\color{codegray}\itshape, breaklines=true, frame=single,
  rulecolor=\color{codegray!30}, showstringspaces=false, tabsize=2,
  morekeywords={library,ggplot,aes,geom_point,geom_col,geom_line,geom_smooth,
    geom_text,geom_label,geom_text_repel,geom_hline,geom_vline,geom_rect,
    annotate,labs,theme_minimal,theme,element_text,coord_flip,
    fct_reorder,scale_fill_manual}
}
\lstdefinestyle{pythonstyle}{
  language=Python, backgroundcolor=\color{codebg}, basicstyle=\ttfamily\scriptsize,
  keywordstyle=\color{codeblue}\bfseries, stringstyle=\color{codegreen},
  commentstyle=\color{codegray}\itshape, breaklines=true, frame=single,
  rulecolor=\color{codegray!30}, showstringspaces=false, tabsize=4,
  morekeywords={import,from,as,pd,plt,sns,np,fig,ax,annotate,text,
    axhline,axvline,axhspan,axvspan,set_title,set_xlabel,set_ylabel}
}

\title{labs(), Annotations \& Storytelling Layers}
\subtitle{Workshop 2 --- Module 8}
\author{Prof.Asoc.\ Endri Raco}
\institute{Data Visualization for Data Scientists \\ UNYT Tirana}
\date{2025/26}

\begin{document}

\begin{frame}
  \titlepage
\end{frame}

\begin{frame}{Module Roadmap}
  \begin{enumerate}
    \item \texttt{labs()}: title, subtitle, caption, axis labels
    \item Title patterns: descriptive, declarative, question
    \item Annotation types: text, arrow, reference line, band, highlight
    \item \texttt{annotate()} vs \texttt{geom\_text()} vs \texttt{geom\_text\_repel()}
    \item Storytelling layers: progressive reveal
    \item Complete annotated chart in R and Python
  \end{enumerate}
  \begin{exampleblock}{Learning Outcome}
    You will transform a data display into a data story by adding
    labels, annotations, and reference marks that guide the reader
    to the intended finding.
  \end{exampleblock}
\end{frame}

\section{labs(): The Label Layer}

\begin{frame}{Anatomy of labs()}
  \begin{center}
    \includegraphics[width=0.82\textwidth]{figures_w02m08/labs_anatomy}
  \end{center}
\end{frame}

\begin{frame}[fragile]{labs() in R vs Python}
  \begin{columns}[T]
    \begin{column}{0.48\textwidth}
      \textbf{R (ggplot2)}
\begin{lstlisting}[style=rstyle]
ggplot(df, aes(x, y, color = g)) +
  geom_point() +
  labs(
    title = "Q4 Revenue Surged 35%",
    subtitle = "Quarterly revenue,
                East region, 2024",
    caption = "Source: Internal DB",
    x = "Displacement (L)",
    y = "Highway MPG",
    color = "Vehicle Class"
  )
\end{lstlisting}
    \end{column}
    \begin{column}{0.48\textwidth}
      \textbf{Python (matplotlib)}
\begin{lstlisting}[style=pythonstyle]
ax.set_title("Q4 Revenue Surged 35%",
    fontsize=14, fontweight="bold")
fig.text(0.13, 0.91,
    "Quarterly revenue, East, 2024",
    fontsize=10, color="grey")
ax.set_xlabel("Displacement (L)")
ax.set_ylabel("Highway MPG")
fig.text(0.13, 0.02,
    "Source: Internal DB",
    fontsize=8, fontstyle="italic",
    color="#888")
ax.legend(title="Vehicle Class")
\end{lstlisting}
    \end{column}
  \end{columns}
  \vspace{4pt}
  \begin{alertblock}{Pro-Tip}
    In ggplot2, \texttt{labs()} sets all labels in one call. In matplotlib,
    title/xlabel/ylabel are on the axes, but subtitle and caption require
    \texttt{fig.text()} with manual coordinates.
  \end{alertblock}
\end{frame}

\section{Title Patterns}

\begin{frame}{Three Title Strategies}
  \begin{center}
    \includegraphics[width=0.92\textwidth]{figures_w02m08/title_patterns}
  \end{center}
  \begin{exampleblock}{Data Insight}
    For executive communication, use \textbf{declarative} titles that
    state the finding as a sentence. For exploratory analysis, use
    \textbf{descriptive} titles. For teaching or interactive contexts,
    use \textbf{question} titles that invite investigation.
  \end{exampleblock}
\end{frame}

\section{Annotation Types}

\begin{frame}{Four Annotation Types in One Chart}
  \begin{center}
    \includegraphics[width=0.88\textwidth]{figures_w02m08/annotation_types}
  \end{center}
\end{frame}

\begin{frame}{Before \& After: Annotations Add Meaning}
  \begin{center}
    \includegraphics[width=0.92\textwidth]{figures_w02m08/before_after_annotations}
  \end{center}
\end{frame}

\begin{frame}[fragile]{Annotations in R (ggplot2)}
\begin{lstlisting}[style=rstyle]
ggplot(df, aes(x = month, y = revenue)) +
  geom_line(color = "#1565C0", linewidth = 1.2) +

  # 1. Text callout with arrow
  annotate("text", x = 8, y = 250, label = "Peak: $250K",
           fontface = "bold", color = "#E53935", size = 4) +
  annotate("segment", x = 8, y = 245, xend = 7, yend = 235,
           arrow = arrow(length = unit(0.2, "cm")), color = "#E53935") +

  # 2. Reference line (mean)
  geom_hline(yintercept = mean(df$revenue),
             linetype = "dashed", color = "grey50") +

  # 3. Reference band (target range)
  annotate("rect", xmin = -Inf, xmax = Inf, ymin = 200, ymax = 220,
           alpha = 0.08, fill = "#2E7D32") +

  # 4. Event marker (vertical line)
  geom_vline(xintercept = 6.5, linetype = "dotted", color = "#E65100") +
  annotate("text", x = 7, y = min(df$revenue), label = "Policy change",
           size = 3, color = "#E65100")
\end{lstlisting}
\end{frame}

\begin{frame}[fragile]{Annotations in Python (matplotlib)}
\begin{lstlisting}[style=pythonstyle]
ax.plot(months, revenue, "-o", color="#1565C0", lw=2)

# 1. Text callout with arrow
ax.annotate("Peak: $250K",
    xy=(7, 235), xytext=(8, 250),            # data coords
    fontsize=9, fontweight="bold", color="#E53935",
    arrowprops=dict(arrowstyle="->", color="#E53935", lw=1.5))

# 2. Reference line (mean)
ax.axhline(y=np.mean(revenue), color="grey", lw=1, linestyle="--")

# 3. Reference band (target range)
ax.axhspan(200, 220, alpha=0.08, color="#2E7D32")

# 4. Event marker (vertical line)
ax.axvline(x=6.5, color="#E65100", lw=0.8, linestyle=":")
ax.text(6.7, min(revenue), "Policy change", fontsize=7, color="#E65100")
\end{lstlisting}
\end{frame}

\section{annotate() vs geom\_text()}

\begin{frame}{Three Labelling Functions}
  \begin{center}
    \includegraphics[width=0.92\textwidth]{figures_w02m08/annotate_vs_geom_text}
  \end{center}
\end{frame}

\begin{frame}[fragile]{geom\_text\_repel: Anti-Overlap Labels}
  \begin{columns}[T]
    \begin{column}{0.48\textwidth}
      \textbf{R}
\begin{lstlisting}[style=rstyle]
library(ggrepel)

ggplot(top10, aes(x = reviews,
    y = rating, label = app)) +
  geom_point(alpha = 0.6) +
  # Labels repel from each other
  geom_text_repel(
    size = 3,
    max.overlaps = 15,
    segment.color = "grey70",
    segment.size = 0.3) +
  scale_x_log10() +
  theme_minimal()
\end{lstlisting}
    \end{column}
    \begin{column}{0.48\textwidth}
      \textbf{Python}
\begin{lstlisting}[style=pythonstyle]
from adjustText import adjust_text

ax.scatter(reviews, rating, s=20)
texts = [ax.text(r, rat, name,
    fontsize=7)
    for r, rat, name in
    zip(reviews, rating, names)]

# Repel overlapping labels
adjust_text(texts,
    arrowprops=dict(
        arrowstyle="-",
        color="grey",
        lw=0.5))
\end{lstlisting}
      {\small Install: \texttt{pip install adjustText}}
    \end{column}
  \end{columns}
\end{frame}

\section{Callout Annotations}

\begin{frame}{Arrow + Text Box: The Callout Pattern}
  \begin{center}
    \includegraphics[width=0.82\textwidth]{figures_w02m08/callout}
  \end{center}
  \begin{alertblock}{Pro-Tip}
    Use callouts sparingly --- one or two per chart maximum. Each callout
    should highlight a \textbf{specific finding}, not just label a data
    point. ``COVID-19 impact'' is a finding; ``2020'' is just a label.
  \end{alertblock}
\end{frame}

\begin{frame}{Reference Lines, Event Markers, Target Bands}
  \begin{center}
    \includegraphics[width=0.92\textwidth]{figures_w02m08/reference_lines}
  \end{center}
\end{frame}

\section{Storytelling Layers}

\begin{frame}{Progressive Reveal: Four Layers of Meaning}
  \begin{center}
    \includegraphics[width=0.95\textwidth]{figures_w02m08/storytelling_layers}
  \end{center}
  \begin{exampleblock}{Data Insight}
    Each layer adds one element of the story: (1) raw data, (2) trend
    context, (3) key events, (4) the finding as a title. This layering
    technique works equally well for presentations (reveal one slide at
    a time) and static reports (the final chart contains all layers).
  \end{exampleblock}
\end{frame}

\section{Complete Examples}

\begin{frame}[fragile]{R: Fully Annotated Bar Chart}
  \begin{columns}[T]
    \begin{column}{0.52\textwidth}
\begin{lstlisting}[style=rstyle]
df |>
  mutate(highlight =
    cat == "SUV") |>
  ggplot(aes(
    x = fct_reorder(cat, n),
    y = n,
    fill = highlight)) +
  geom_col(width = 0.5,
    show.legend = FALSE) +
  geom_text(aes(label = n),
    hjust = -0.2, size = 3.5,
    fontface = "bold") +
  scale_fill_manual(
    values = c("TRUE" = "#E53935",
               "FALSE" = "#BBBBBB")) +
  geom_hline(yintercept = mean_n,
    linetype = "dashed",
    color = "grey50") +
  coord_flip() +
  labs(
    title = "SUV Leads with 62",
    subtitle = "Vehicle count, 2024",
    caption = "Source: mpg | UNYT",
    x = NULL, y = NULL) +
  theme_minimal()
\end{lstlisting}
    \end{column}
    \begin{column}{0.45\textwidth}
      \includegraphics[width=\textwidth]{figures_w02m08/r_annotated}

      \vspace{4pt}
      {\small Annotation checklist:
      \begin{enumerate}
        \item Declarative title
        \item Contextual subtitle
        \item Source caption
        \item Direct bar labels
        \item Grey-accent colour
        \item Mean reference line
      \end{enumerate}}
    \end{column}
  \end{columns}
\end{frame}

\begin{frame}[fragile]{Python: Same Annotated Chart}
  \begin{columns}[T]
    \begin{column}{0.52\textwidth}
\begin{lstlisting}[style=pythonstyle]
fig, ax = plt.subplots(figsize=(6,5))

colors = ["#BBBBBB"] * 5
colors[-1] = "#2E7D32"  # accent

ax.barh(sorted_cats, sorted_vals,
    color=colors, height=0.5)

# Direct labels
for i, v in enumerate(sorted_vals):
    ax.text(v + 0.8, i, str(v),
        va="center", fontsize=9,
        fontweight="bold",
        color="#2E7D32"
            if i == 4 else "#888")

# Reference line
ax.axvline(x=np.mean(sorted_vals),
    color="#888", lw=0.8, ls="--")

# Title + subtitle + caption
ax.set_title("SUV Leads with 62",
    fontsize=11, fontweight="bold")
fig.text(0.13, 0.88,
    "Vehicle count, 2024",
    fontsize=9, color="grey")
fig.text(0.13, 0.02,
    "Source: mpg | UNYT",
    fontsize=7, fontstyle="italic")
\end{lstlisting}
    \end{column}
    \begin{column}{0.45\textwidth}
      \includegraphics[width=\textwidth]{figures_w02m08/py_annotated}
    \end{column}
  \end{columns}
\end{frame}

\section{Synthesis}

\begin{frame}{Annotation Checklist}
  Before finalising any chart, verify:
  \begin{enumerate}
    \item \textbf{Title}: Does it state the finding (declarative) or
          describe the chart (descriptive)?
    \item \textbf{Subtitle}: Does it add context (time period, filter,
          data source)?
    \item \textbf{Caption}: Does it credit the source and note methodology?
    \item \textbf{Axis labels}: Do they include units (\$, \%, kg)?
    \item \textbf{Direct labels}: Are key values shown on the chart
          (eliminating need to read axes)?
    \item \textbf{Reference marks}: Is there a mean line, target band,
          or event marker for context?
    \item \textbf{Accent colour}: Is the key finding visually highlighted?
  \end{enumerate}
\end{frame}

\begin{frame}{Key Takeaways}
  \begin{enumerate}
    \item \textbf{labs()} sets title, subtitle, caption, and axis labels.
          Use declarative titles for executive audiences.
    \item Four annotation types: \textbf{text callout} (arrow + label),
          \textbf{reference line} (mean, target), \textbf{band} (range),
          \textbf{highlight region} (event period).
    \item \textbf{annotate()} for fixed-position single annotations;
          \textbf{geom\_text()} for data-driven labels;
          \textbf{geom\_text\_repel()} for anti-overlap.
    \item \textbf{Storytelling layers} build meaning progressively:
          data $\rightarrow$ trend $\rightarrow$ events $\rightarrow$ finding.
    \item Every production chart needs: title, subtitle, caption, units,
          and at least one annotation that guides the reader.
  \end{enumerate}
\end{frame}

\begin{frame}{Essential Readings}
  \begin{itemize}
    \item Knaflic, C.~N. (2015). \emph{Storytelling with Data},
          Chapter 4 (Focus Your Audience's Attention).
    \item Schwabish, J. (2021). \emph{Better Data Visualizations},
          Chapter 5 (Annotations).
    \item Wickham, H. (2016). \emph{ggplot2}, Chapter 8.3 (Annotations).
    \item Wilke, C.~O. (2019). \emph{Fundamentals of Data Visualization},
          Chapter 22 (Titles, Captions, and Tables).
  \end{itemize}
  \vspace{6pt}
  \textbf{Next Module}: Tidy Data \& Data Wrangling for Viz (W02-M09)
\end{frame}

\end{document}
